%% LyX 2.4.4 created this file.  For more info, see https://www.lyx.org/.
%% Do not edit unless you really know what you are doing.
\documentclass[oneside,english]{book}
\usepackage[T1]{fontenc}
\usepackage[latin9]{inputenc}
\setcounter{secnumdepth}{3}
\setcounter{tocdepth}{3}
\usepackage{amsmath}
\usepackage{amsthm}
\usepackage{amssymb}
\usepackage{cancel}
\usepackage{geometry}
\geometry{verbose,tmargin=1cm,bmargin=1cm,lmargin=1.5cm,rmargin=1.5cm}

\makeatletter
%%%%%%%%%%%%%%%%%%%%%%%%%%%%%% Textclass specific LaTeX commands.
\theoremstyle{definition}
\newtheorem{defn}{\protect\definitionname}
\theoremstyle{definition}
\newtheorem{example}{\protect\examplename}
\theoremstyle{plain}
\newtheorem{prop}{\protect\propositionname}
\ifx\proof\undefined\
  \newenvironment{proof}[1][\proofname]{\par
    \normalfont\topsep6\p@\@plus6\p@\relax
    \trivlist
    \itemindent\parindent
    \item[\hskip\labelsep
          \scshape
      #1]\ignorespaces
  }{%
    \endtrivlist\@endpefalse
  }
  \providecommand{\proofname}{Proof}
\fi
\theoremstyle{definition}
\newtheorem{xca}{\protect\exercisename}
\theoremstyle{plain}
\newtheorem{thm}{\protect\theoremname}

%%%%%%%%%%%%%%%%%%%%%%%%%%%%%% User specified LaTeX commands.
\usepackage{hyperref}
\usepackage[usenames,dvipsnames]{xcolor} %adds additional color options
\hypersetup{colorlinks=true, urlcolor=Blue}%Blue

\usepackage{multicol}

\usepackage{tikz}
\usepackage{tikz-3dplot}
\usetikzlibrary{calc,arrows,shapes,backgrounds,matrix,shapes.geometric,fit,trees,positioning}

\usetikzlibrary{patterns}
\usetikzlibrary{plotmarks}
\usepackage{pgfplots}
\pgfplotsset{compat=newest}

\usepackage{calc}
\usepackage[nomessages]{fp}

\usepackage{datetime}

%%%%%commands
\newcommand{\mb}[1]{$\mathbb{#1}$}
\newcommand{\funct}[2]{$f\colon#1\rightarrow#2$}
% Added by lyx2lyx
\setlength{\parskip}{\smallskipamount}
\setlength{\parindent}{0pt}

\makeatother

\usepackage{babel}
\providecommand{\definitionname}{Definition}
\providecommand{\examplename}{Example}
\providecommand{\exercisename}{Exercise}
\providecommand{\propositionname}{Proposition}
\providecommand{\theoremname}{Theorem}

\begin{document}

\chapter{\label{chap:Rings}Rings}

\section{Rings and Subrings}

We now continue to generalize the familiar arithmetic operations of
addition and multiplication. Just as with groups, results from the
algebra of arithmetic will be extended to a host of non-numerical
objects, e.g., matrices, functions, polynomials, etc.
\begin{itemize}
\item From here on out we will use the following notations unless stated
otherwise:

\begin{itemize}
\item []$+$ for the ``addition'' operation, e.g., $a+b$ (instead of
$\oplus$).
\item []$\cdot$ or juxtaposition for the ``multiplication'' operation,
e.g., $a\cdot b$ or $ab$ (instead of $\otimes$).
\item []''$-a$'' will denote the additive inverse of an element $a$,
e.g., $-a+a=0$ (instead of ``$a^{-1}$'').
\item []''$a^{-1}$'' will denote the multiplicative inverse of an element
$a$, e.g., $a^{-1}a=1$.
\end{itemize}
\end{itemize}
\begin{defn}
A \emph{ring} is a nonempty set, say $R$, with two binary operations
(usually written as multiplication and addition)\footnote{\emph{But of course the operations are not necessarily addition and
multiplication. $+$ and $\times$; $\oplus$ and $\otimes$; are
among the different notations that may be used in varying texts. }} that satisfies the following:

\begin{enumerate}
\item If $a,b\in R$ then $a+b\in R$. \hfill(closed under addition)
\item $a+(b+c)=(a+b)+c$ \hfill(associative addition)
\item $a+b=b+a$ \hfill(commutative addition)
\item There exists an $0_{R}$ such that $a+0_{R}=0_{R}+a=a$ \hfill(additive
identity)
\item There exists an $-a\in R$ such that $a+(-a)=0_{R}$\\
for every $a\in R$. \hfill(additive inverse)
\item If $a,b\in R$ then $ab\in R$. \hfill(closed under multiplication)
\item $a(bc)=(ab)c$ \hfill(associative multiplication)
\item $a(b+c)=ab+ac$ and $(a+b)c=ac+bc$ \hfill(distribution holds).
\end{enumerate}
\end{defn}
\begin{itemize}
\item Note that parts 1-5 say that a ring is an \emph{abelian group.}\footnote{Commutative groups are also called abelian in honor of Norwegian mathematician
Niels Abel.}\emph{ }
\end{itemize}
\begin{example}
The set of integers, $\mathbb{Z}$, under addition and multiplication
is a ring.
\end{example}

\begin{example}
$\mathbb{Z}_{4}=\{0,1,2,3\}$ ($\mathbb{Z}$ modulo $4$) is also
a ring under addition and multiplication as in $\mathbb{Z}_{n}$ for
any $n\in\mathbb{N}$.
\end{example}

\begin{example}
The set of $2$ by $2$ matrices with real numbers as entries,\\
 $\mathcal{M}_{2}(\mathbb{R})=\left\{ \begin{bmatrix}\begin{array}{rr}
a & b\\
c & d
\end{array}\end{bmatrix}\colon a,b,c,d\in\mathbb{R}\right\} $.
\end{example}
\begin{itemize}
\item []Closure, associativity, and distribution can be shown using the
usual rules of matrix multiplication and addition. $\mathcal{M}_{2}(\mathbb{R})$
is commutative under addition (but not multiplication), the additive
inverse is $\begin{bmatrix}\begin{array}{rr}
-a & -b\\
-c & -d
\end{array}\end{bmatrix}$, the additive identity is $\begin{bmatrix}\begin{array}{rr}
0 & 0\\
0 & 0
\end{array}\end{bmatrix}$, and the multiplicative identity is $\begin{bmatrix}\begin{array}{rr}
1 & 0\\
0 & 1
\end{array}\end{bmatrix}$.
\end{itemize}

\begin{prop}
If $a$ and \textbf{$b$} are two elements in $R$ then $(-a)(-b)=ab$.
\end{prop}
\begin{proof}
By definition $a+(-a)=0$. $(a+(-a))(-b)=0$ implies $a(-b)+(-a)(-b)=0$
by distribution. The additive inverses of $a(-b)$ is unique thus
$a(-b)+(-a)(-b)=0$ implies that $(-a)(-b)=ab$.
\end{proof}
\begin{itemize}
\item In particular for rings like $\mathbb{Z}$ and $\mathbb{R}$, the
product of two negatives is a positive. Note the importance of distribution.
\end{itemize}
\begin{xca}
Show that $\mathbb{Z}_{6}$ is a ring with unity under addition and
multiplication.
\end{xca}

\begin{xca}
Show that the set $\sqrt{2}\mathbb{Z}=\{x\colon x=m+\sqrt{2}n\mbox{ for some \ensuremath{m,n\in\mathbb{Z}}}\}$
is a ring.
\end{xca}

\begin{xca}
Let $R$ and $S$ be rings. Define addition and multiplication on
the Cartesian product $R\times S$ by,
\[
(r,s)+(r',s')=(r+r',s+s')\mbox{ and }(r,s)(r',s')=(rr',ss')
\]
for $r,r'\in R$ and $s,s'\in S$. Prove that $R\times S$ is a ring. 
\end{xca}
\begin{defn}
A \emph{commutative ring }is a ring $R$ such that $ab=ba$ for all
$a,b\in R$.
\end{defn}

\begin{defn}
A ring with \emph{identity}\footnote{The term \emph{unity} is often used in place identity. A ring with
an identity element is said to ``have unity''.}\emph{ }is a ring $R$ with an element $1_{R}$ such that $a1_{R}=1_{R}a$
for all $a\in R$.
\end{defn}
\begin{itemize}
\item It is important to check that a potential $1_{R}$ is both a \emph{right
identity, $a1_{R}=a$}; and a \emph{left identity}, $1_{R}b=b$.
\end{itemize}
\begin{example}
$\mathbb{C}$, $\mathbb{R}$, $\mathbb{Q}$, and $\mathbb{Z}$ are
all commutative rings with identity under the usual addition and multiplication.
\end{example}

\begin{example}
\label{ex-Zn-is-comm-ring}$\mathbb{Z}_{n}=\{0,1,2,\dots,n-1\}$ where
$n\geq2$ is a commutative ring with identity under the usual addition
and multiplication (recall that $n$ is called the \emph{modulus}).
\end{example}
\begin{itemize}
\item To add $x+y\mod n$ find $(x+y)\div n$ and find the remainder, say
$r$, then $(x+y)\mod n=r$. In the same way we can find $(xy)\mod n$.
For example, $(3\cdot7)\mod 8=5$ since $2\cdot8+5=3\cdot7$ (divide
by $8$ and take the remainder). 

\begin{proof}
Since $r<n$, $\mathbb{Z}_{n}$ will be closed under addition and
multiplication. Parts 2, 3, 4, and 7 hold because of the normal rules
of arithmetic. Which also give multiplicative commutativity. $0$
and $1$ are the additive and multiplicative identities respectively.

Which only leaves distribution, part 8, and additive inverse, part
5. Let $a\in\mathbb{Z}_{n}$ then $a<n$ so then $n-a\in\mathbb{Z}_{n}$,
and $a+(n-a)=0$. For distribution, let $a,b,c\in\mathbb{Z}_{n}$.
$a(b+c)\mod n=(ab+ac)\mod n=ab\mod n+ac\mod n$ (this was shown last
semester.
\end{proof}
\end{itemize}

\begin{example}
$E=\{x\colon x=2n\mbox{ for some }n\in\mathbb{Z}\}$, the set of even
integers is a commutative ring. However it does \emph{not }have identity. 
\end{example}
\begin{itemize}
\item Keep in mind that a ring is \emph{not }a group under multiplication.
Elements in a ring do not necessarily have a multiplicative inverse,
and rings do not necessarily have a multiplicative identity. Keep
examples of rings in mind which do not have these properties. 
\end{itemize}
\begin{xca}
Prove that a ring that is cyclic\footnote{Recall that a group $G$ is \emph{cyclic} if for every $x\in G$ there
is some $g\in G$ such that $x=g+g+\cdots+g$. In which case we write
$G=\left\langle g\right\rangle $ and say that $G$ is \emph{generated
}by $g$, e.g., $\mathbb{Z}_{3}=\left\langle 1\right\rangle =\left\langle 2\right\rangle $
since $2+2=1$, $2=2$, and $2+2+2=0$ $\mod 3$.} under addition and commutative (under multiplication).
\end{xca}
\begin{thm}
If a ring has an identity, it is unique. 
\end{thm}
\begin{proof}
Suppose to the contrary that a ring $R$ has two identity elements
$1$ and $1'$. $a1=a$ for all $a\in R$ including $1'\in R$. Thus
$1'\cdot1=1$' and $1\cdot1'=1$ . By the definition of a ring $1'=1\cdot1'=1'\cdot1=1$.
Therefore $1'=1$, and the multiplicative identity is unique.
\end{proof}
\begin{xca}
\label{ex-funct-are-ring}Let $F$ be the set of all functions from
$\mathbb{R}$ to $\mathbb{R}$. Define addition and multiplication
as 
\[
(f+g)(x)=f(x)+g(x)\mbox{ and }(fg)(x)=f(x)g(x).
\]
Show that $F$ is a commutative ring with identity.
\end{xca}

\begin{xca}
If a ring element has a multiplicative inverse, it is unique.
\end{xca}

\section{Subrings}
\begin{defn}
Let $R$ be a ring and $S$ be a subset of $R$.\footnote{Notated: $S\subseteq R$.}
If $S$ is a ring under the same operations as $R$, then $S$ is
a \emph{subring}.
\end{defn}
\begin{example}
$\mathbb{Q}$ is a subring of $\mathbb{R}$, and $\mathbb{Z}$ is
a subring of $\mathbb{Q}$ and $\mathbb{R}$.
\end{example}

\begin{example}
The set of even integers is a subring of $\mathbb{Z}$, but the set
of odd integers is not since it is not closed under addition. 
\end{example}
The following two theorems provide shortcuts when trying to show that
a subset is a subring. Particularity, The second theorem can save
quite a bit of work.
\begin{thm}
Suppose that $R$ is a ring and $S\subseteq R$ such that the following
hold:

\begin{enumerate}
\item If $a,b\in S$ then $a+b\in S$ (closed under addition)
\item If $a,b\in S$ then $ab\in S$ (closed under multiplication)
\item $0_{R}\in S$.
\item If $a\in S$ then there is an $-a\in S$ such that $a+(-a)=0_{R}$.
\end{enumerate}
Then $S$ is a subring of $R$.
\end{thm}
\begin{proof}
Parts 2, 3, 7, and 8 of the Ring definition hold for all elements
in $R$ since it is a ring. So they hold for all elements in $S$.
The remaining requirements 1-4 in the theorem satisfy the remaining
requirements. 
\end{proof}
\begin{itemize}
\item Rings and thus subrings must be nonempty. A requirement which is covered
by 3. Since $0_{R}\in S\Rightarrow$ $S\neq\emptyset$. 
\end{itemize}
\begin{thm}
\label{thm: subring test}A nonempty subset $S$ of a ring $R$ is
a subring if whenever $a,b\in S$:

\begin{enumerate}
\item $ab\in S$ (closed under multiplication)
\item $a-b\in S$ (closed under subtraction)
\end{enumerate}
\end{thm}
\begin{proof}
If $a-b\in S$ for all $a,b\in S$ then $S$ is a subgroup of $R$
(subgroup test), and since $R$ is additively commutative it follows
that $S$ is a abelian subgroup of $R$ under addition (this also
means that $0_{R}$). $S$ is multiplicatively and additively associative;
and $S$ is distributive over addition because these properties hold
in $R$ .
\end{proof}
\begin{xca}
Show that the set $3\mathbb{Z}=\{x\colon x=3n\mbox{ for some \ensuremath{n\in\mathbb{Z}}}\}$
is a subring of $\mathbb{Z}$.
\end{xca}

\begin{xca}
Show that the set $k\mathbb{Z}=\{x\colon x=kn\mbox{ for some \ensuremath{n\in\mathbb{Z}}}\}$
for any integer $k$ is a subring of $\mathbb{Z}$.
\end{xca}

\begin{xca}
Prove that the intersection of any two subrings of a ring $R$ is
also a subring of $R$.
\end{xca}

\end{document}
